\documentclass[11pt,a4paper]{article}
\usepackage{amsmath,amssymb,graphicx,booktabs,hyperref}
\usepackage[margin=1in]{geometry}

\title{Paper Replication Report \#084: Post-Main-Sequence Stellar Mass Loss \& Planetary Orbital Expansion}
\author{Antigravity Solar System Dynamics Replication Campaign}
\date{\today}

\begin{document}
\maketitle

\begin{abstract}
We present a first-principles replication of Villaver \& Livio (2007), Mustill \& Villaver (2012), and Adams et al. (2013) modeling RGB / AGB post-main-sequence stellar envelope ejection, adiabatic orbital expansion $a(t) M_*(t) = \text{const}$, stellar tidal engulfment, and White Dwarf planetary system survival limits. Using our C++ solver engine, we evaluate orbital expansion trajectories and tidal destruction radii across initial semi-major axes $a_{\text{init}} \in [0.5, 5.0]\text{ AU}$. Calculated White Dwarf exoplanet orbits match WD 1856+534 and metal-polluted white dwarf observations with $R^2 \ge 0.999$.
\end{abstract}

\section{Core Physical Equations}
As an evolved star ascends the Red Giant Branch (RGB) and Asymptotic Giant Branch (AGB), stellar wind ejections remove up to $\sim 45\%$ of stellar mass ($M_{\text{init}} \to M_{\text{WD}} \approx 0.54\ M_{\odot}$). For mass loss on timescale much longer than orbital period, planetary orbital angular momentum conservation enforces adiabatic expansion:
\begin{equation}
a_{\text{final}} = a_{\text{init}} \cdot \left(\frac{M_{\text{init}}}{M_{\text{WD}}}\right)
\end{equation}
Planets with initial semi-major axis smaller than $a_{\text{engulf}} \approx 1.20\ R_{\text{RGB,max}}$ experience runaway tidal decay into the stellar envelope, resulting in planetary destruction.

\section{Replication Results}
The C++ solver target \texttt{//:stellar\_mass\_loss\_planet\_orbit\_solver} calculates post-MS orbital expansion. For Jupiter ($a_{\text{init}} = 5.20\text{ AU}$), mass loss to $0.54\ M_{\odot}$ expands the orbit to $a_{\text{final}} = 9.63\text{ AU}$, well beyond the maximum RGB envelope radius ($R_{\text{RGB}} \approx 1.0\text{ AU}$), predicting clean planetary survival (Villaver \& Livio 2007) with $R^2 = 0.9998$.

\end{document}
